% ============================================================================
% Chapter 6: Multi-Party Document Signing, Privacy, and Legal Compliance
% Study Guide for Cardano Credential Verification
% ============================================================================

\chapter{Multi-Party Document Signing, Privacy, and Legal Compliance}
\label{ch:signing-privacy-legal}

The preceding chapters established how a single professional obtains a license
NFT, holds signature and validity tokens, and participates in the credential
ecosystem.  This chapter confronts three harder questions that arise the moment
the system moves from individual credentials to real-world workflows:
\emph{How do multiple professionals sign a single document while proving their
credentials?}  \emph{How does a public blockchain protect private information?}
\emph{Does any of this hold up in court?}

These three concerns---multi-party signing, privacy architecture, and legal
compliance---are deeply intertwined.  A signing workflow that leaks personal
data violates GDPR.  A privacy scheme that prevents verification defeats the
legal admissibility of the signature.  A legally compliant system that cannot
handle concurrent signers is useless for the building-design reviews,
multi-attorney closings, and medical peer-review panels that constitute the
system's primary use cases.  This chapter addresses all three in sequence and
then shows how they compose into a coherent whole.

% Load TikZ libraries needed for this chapter
\tikzset{
  guidebox/.style={draw, rounded corners=4pt, thick, align=center,
    font=\sffamily\small, fill=white, minimum height=1cm},
  guidearrow/.style={-{Latex[length=2.5mm]}, thick, color=gray!70!black},
  guidelabel/.style={font=\sffamily\footnotesize, text=gray!50!black},
}


% ============================================================================
% PART A: MULTI-PARTY DOCUMENT SIGNING
% ============================================================================

\vspace{0.5cm}
\noindent\rule{\textwidth}{1.2pt}
\begin{center}
  {\Large\sffamily\bfseries Part A: Multi-Party Document Signing}
\end{center}
\noindent\rule{\textwidth}{1.2pt}
\vspace{0.3cm}


% ----------------------------------------------------------------------------
\section{The Signing Problem}
\label{sec:signing-problem}

Consider a concrete scenario.  A seismic retrofit design for a hospital in
Nevada requires signatures from two licensed Professional Engineers (one
structural, one geotechnical) and one licensed notary public.  Each
professional must \emph{prove} that they hold a valid credential \emph{at the
exact moment} they sign.  If either engineer's license was suspended yesterday,
the signature is worthless---and potentially fraudulent.

Under current practice, this verification happens through one of three
mechanisms, all of which are inadequate:

\begin{itemize}[leftmargin=*, nosep, itemsep=4pt]
  \item \textbf{Wet signatures with manual verification.}  Each professional
    signs a physical document.  The reviewing authority then contacts each
    licensing board---potentially across multiple states---to confirm that each
    license was active on the date of signing.  This process takes days to
    weeks and produces no machine-verifiable audit trail.
  \item \textbf{Disconnected e-signature tools.}  Services like DocuSign or
    Adobe Sign capture digital signatures, but they verify \emph{identity}
    (``this email address clicked `sign'\,''), not \emph{credentials} (``this
    person held an active PE license in California at 14:32 UTC on February
    19, 2026'').  The credential verification step is still a separate,
    manual process.
  \item \textbf{PDF digital signatures with PKI.}  X.509 certificate-based
    PDF signatures provide cryptographic proof of identity, but the
    certificate's validity period is coarse (typically one year), the
    revocation check depends on the issuing CA's CRL or OCSP responder
    being online, and there is no standard mechanism for binding the
    signature to a \emph{professional license} rather than a \emph{personal
    identity}.
\end{itemize}

The fundamental gap is that no existing commercial signing tool answers the
question: ``Was this signer a licensed professional in good standing at the
moment of signing?''  The credential verification system described in this
guide closes that gap by making credential status an \emph{on-chain,
machine-verifiable precondition} of the signing act itself.

The multi-party dimension compounds the difficulty.  It is not enough for one
signer to prove credentials; every required signer must do so, and the final
work product must atomically record that \emph{all} required signatures were
collected with \emph{all} credentials verified.  Partial completion---two of
three signatures---must not be treated as final.  Yet the system must also
handle the practical reality that three professionals in different time zones
will not sign simultaneously.


% ----------------------------------------------------------------------------
\section{The Six-Phase Signing Workflow}
\label{sec:six-phases}

The multi-party signing workflow proceeds through six distinct phases, each
with specific on-chain and off-chain operations.  Understanding this workflow
in detail is essential because it is the mechanism by which the credential
system transforms from a \emph{credential registry} into a \emph{document
signing infrastructure} with built-in credential verification.

\subsection{Phase 1: Setup}

The process begins when an authority (or an authorized coordinator) creates a
\textbf{work product entry} in the system.  This entry records three pieces of
information:

\begin{enumerate}[nosep]
  \item \textbf{Document hash.}  The SHA-256 hash of the document to be
    signed.  This 32-byte value uniquely identifies the document and ensures
    that any modification---even changing a single bit---produces a
    completely different hash.  The document itself is never stored on-chain;
    only its fingerprint is.
  \item \textbf{Required signers.}  A list of public key hashes (PKHs)
    identifying the professionals who must sign.  Each PKH corresponds to
    a Cardano wallet address derived from the signer's CIP-1852 HD wallet.
  \item \textbf{Work product metadata.}  A human-readable description,
    project identifier, and creation timestamp.  This metadata is stored in
    the off-chain database and referenced by the on-chain records.
\end{enumerate}

At this stage, nothing is written to the blockchain.  The work product entry
exists only in the local SQLite database.  The purpose of Phase~1 is to
establish the \emph{requirements}---what document, which signers---before
any on-chain resources are committed.

\subsection{Phase 2: Credential Verification}

Before any signer can deposit tokens, the system verifies that each required
signer holds the necessary credentials.  For each signer, three checks are
performed:

\begin{enumerate}[nosep]
  \item \textbf{License NFT (CIP-68).}  The system queries the Cardano
    blockchain to confirm that the signer's wallet contains a license NFT
    minted under a recognized authority's policy ID.  Critically, it also
    reads the CIP-68 \emph{Reference Token} (asset label 100) held by the
    authority to confirm that the inline datum field \texttt{status} is set
    to \texttt{active}---not \texttt{revoked} or \texttt{suspended}.  This
    is the mechanism by which real-time credential status is verified: even
    if the signer possesses the User Token (asset label 222), the authority
    can unilaterally revoke the credential by updating the Reference Token's
    datum.
  \item \textbf{Signature tokens.}  The signer must hold at least one
    signature token (a fungible native token) minted under the same
    authority's policy.  Each signing act consumes one signature token, so
    the system checks the signer's wallet balance.
  \item \textbf{Validity token.}  The signer must hold a validity token
    whose expiry slot is in the \emph{future}.  The validity token represents
    current good standing---it may be linked to dues payments, continuing
    education, or periodic reauthorization.  If the validity token has
    expired, the signer must renew it (typically by paying dues through the
    Dues Enforcement Contract) before proceeding.
\end{enumerate}

If any signer fails any check, the system halts and reports which signer
failed which requirement.  This prevents wasted on-chain transactions and
provides clear feedback for remediation.

\subsection{Phase 3: Signing (Token Deposit)}

This is the critical on-chain phase.  Each signer constructs and submits a
transaction that deposits \textbf{one signature token} and \textbf{one
validity token} to the \texttt{SignatureCollectionValidator} script address.
The transaction also attaches a \texttt{SignerDatum} (detailed in
Section~\ref{sec:signer-datum}) recording the signer's identity, the
document hash, and temporal information.

The crucial design decision is that \textbf{each signer creates an
independent UTxO} at the script address.  This means that if three signers
deposit, there are three separate UTxOs sitting at the validator address,
each containing one signer's tokens and datum.  This is not an incidental
implementation choice---it is the solution to the eUTxO concurrency
problem (Section~\ref{sec:eutxo-concurrency}).

Each signer's deposit transaction uses the \textsc{Collect} redeemer
(redeemer index 0), which tells the validator: ``I am a required signer,
and I am depositing my tokens for this work product.''  The validator script
checks that the deposited tokens come from a recognized policy ID and that
the datum fields are well-formed.

\subsection{Phase 4: Collection}

As deposits arrive, the system tracks the state of the work product: which
signers have deposited and which have not.  This tracking occurs both
off-chain (in the SQLite database, for dashboard display and status queries)
and on-chain (by querying the validator address for UTxOs whose datums
reference the work product's document hash).

The collection phase is inherently asynchronous.  Signer~A might deposit at
9:00~AM EST, Signer~B at 2:00~PM EST, and Signer~C at 11:00~PM EST the
following day.  The system accommodates this by not requiring any ordering
or timing constraints on deposits---only that all required deposits are
present before finalization.

During this phase, any signer who has already deposited may invoke the
\textsc{Reclaim} redeemer (redeemer index 2) to cancel their participation
and recover their tokens.  This is important for practical workflow
management: if a project is cancelled or a signer needs to be replaced,
deposited tokens are not permanently locked.

\subsection{Phase 5: Finalization}

Once all required signers have deposited, the authority submits a
\textsc{Finalize} transaction (redeemer index 1).  This is the most
consequential on-chain operation in the entire workflow because it
\textbf{atomically consumes ALL per-signer UTxOs} in a single transaction.

The atomicity guarantee is critical.  Either \emph{all} deposits are
consumed and the work product is finalized, or \emph{none} are---there is
no intermediate state where some signatures are finalized and others are
not.  This is a natural property of Cardano's transaction model: a
transaction either succeeds completely or fails completely.

The finalize transaction:
\begin{itemize}[nosep]
  \item Consumes each signer's deposit UTxO (one input per signer)
  \item Verifies that the authority's signature is present on the transaction
  \item Confirms that the set of consumed UTxOs covers all required signers
  \item Records the finalization on-chain as a completed transaction with
    all signer datums preserved in the transaction metadata
\end{itemize}

After finalization, the signature tokens and validity tokens are consumed
(effectively burned from the signers' perspective, as they now reside in
the spent transaction history rather than in any wallet).

\subsection{Phase 6: Verification}

After finalization, \emph{any party}---a building inspector, a court, a
regulatory body, a counterparty---can verify the signed work product by
querying the blockchain.  The verification process confirms:

\begin{itemize}[nosep]
  \item \textbf{Who signed:}  Each signer's PKH is recorded in the
    \texttt{SignerDatum} attached to their deposit UTxO.
  \item \textbf{When they signed:}  The deposit slot timestamp records
    when each signer's deposit transaction was included in a block.
  \item \textbf{What they signed:}  The SHA-256 document hash in each
    datum binds the signature to a specific document version.
  \item \textbf{Credential status at signing time:}  The validity token's
    expiry slot, combined with the deposit slot, proves that the signer's
    credentials were valid at the moment of signing.
  \item \textbf{Finalization:}  The finalize transaction's existence proves
    that the authority confirmed all signatures were collected and valid.
\end{itemize}

This verification requires only a Cardano blockchain query---no contact with
any licensing board, no reliance on any centralized server, no trust in any
single party's attestation.

\vspace{0.5cm}
\begin{figure}[H]
\centering
\begin{tikzpicture}[scale=0.82, every node/.style={transform shape}]
  % Actors
  \node[guidebox, fill=gray!10, minimum width=2cm] (authority) at (0,0)
    {Authority};
  \node[guidebox, fill=blue!8, minimum width=2cm, text width=2cm] (signerA) at (4,0)
    {Signer A (PE Structural)};
  \node[guidebox, fill=green!8, minimum width=2cm, text width=2cm] (signerB) at (8,0)
    {Signer B (PE Geotech)};
  \node[guidebox, fill=orange!8, minimum width=2cm, text width=2cm] (signerC) at (12,0)
    {Signer C (Notary)};

  % Lifelines
  \draw[dashed, gray!40] (0,-0.6) -- (0,-18.5);
  \draw[dashed, gray!40] (4,-0.6) -- (4,-18.5);
  \draw[dashed, gray!40] (8,-0.6) -- (8,-18.5);
  \draw[dashed, gray!40] (12,-0.6) -- (12,-18.5);

  % Phase 1: Setup
  \draw[thick, blue!60] (-1.8,-1.2) -- (13.5,-1.2);
  \node[guidelabel, anchor=west, font=\sffamily\footnotesize\bfseries,
    text=blue!70] at (-1.8,-1.5) {Phase 1: Setup};
  \draw[guidearrow] (0,-2) -- node[above, font=\sffamily\tiny]
    {Create work product} node[below, font=\sffamily\tiny]
    {(doc hash + signer list)} (0,-2);
  \node[guidebox, fill=blue!5, minimum width=2.2cm, minimum height=0.6cm,
    font=\sffamily\tiny, text width=2cm] at (0,-2.4) {SHA-256 hash + 3 PKHs};

  % Phase 2: Verify
  \draw[thick, green!60] (-1.8,-3.2) -- (13.5,-3.2);
  \node[guidelabel, anchor=west, font=\sffamily\footnotesize\bfseries,
    text=green!70!black] at (-1.8,-3.5) {Phase 2: Verify Credentials};
  \draw[guidearrow] (0,-4) -- node[above, font=\sffamily\tiny]
    {Check: NFT + sig tokens + validity} (4,-4);
  \draw[guidearrow] (0,-4.5) -- node[above, font=\sffamily\tiny]
    {Check: NFT + sig tokens + validity} (8,-4.5);
  \draw[guidearrow] (0,-5) -- node[above, font=\sffamily\tiny]
    {Check: NFT + sig tokens + validity} (12,-5);
  \node[guidebox, fill=green!5, minimum width=1cm, minimum height=0.4cm,
    font=\sffamily\tiny] at (4,-4.4) {\checkmark};
  \node[guidebox, fill=green!5, minimum width=1cm, minimum height=0.4cm,
    font=\sffamily\tiny] at (8,-4.9) {\checkmark};
  \node[guidebox, fill=green!5, minimum width=1cm, minimum height=0.4cm,
    font=\sffamily\tiny] at (12,-5.4) {\checkmark};

  % Phase 3: Sign
  \draw[thick, yellow!70!black] (-1.8,-6.0) -- (13.5,-6.0);
  \node[guidelabel, anchor=west, font=\sffamily\footnotesize\bfseries,
    text=yellow!60!black] at (-1.8,-6.3) {Phase 3: Sign (Deposit Tokens)};
  \draw[guidearrow, blue!70] (4,-6.8) -- node[above, font=\sffamily\tiny]
    {Deposit 1 sig + 1 val token} node[below, font=\sffamily\tiny]
    {(independent UTxO \#1)} (1,-6.8);
  \draw[guidearrow, green!70!black] (8,-7.6) -- node[above, font=\sffamily\tiny]
    {Deposit 1 sig + 1 val token} node[below, font=\sffamily\tiny]
    {(independent UTxO \#2)} (1,-7.6);
  \draw[guidearrow, orange!70!black] (12,-8.4) -- node[above, font=\sffamily\tiny]
    {Deposit 1 sig + 1 val token} node[below, font=\sffamily\tiny]
    {(independent UTxO \#3)} (1,-8.4);

  % Validator box
  \node[guidebox, fill=purple!8, minimum width=2.5cm, minimum height=1.2cm,
    font=\sffamily\tiny, text width=2.2cm] (valbox) at (-1.2,-9.6)
    {Validator Address\\[2pt]UTxO \#1 (A)\\UTxO \#2 (B)\\UTxO \#3 (C)};

  % Phase 4: Collect
  \draw[thick, orange!60] (-1.8,-10.8) -- (13.5,-10.8);
  \node[guidelabel, anchor=west, font=\sffamily\footnotesize\bfseries,
    text=orange!70!black] at (-1.8,-11.1) {Phase 4: Collection (Track Status)};
  \node[guidebox, fill=orange!5, minimum width=3cm, minimum height=0.5cm,
    font=\sffamily\tiny] at (6,-11.6) {3/3 deposits collected -- ready to finalize};

  % Phase 5: Finalize
  \draw[thick, red!60] (-1.8,-12.3) -- (13.5,-12.3);
  \node[guidelabel, anchor=west, font=\sffamily\footnotesize\bfseries,
    text=red!70!black] at (-1.8,-12.6) {Phase 5: Finalize (Atomic Sweep)};
  \draw[guidearrow, red!60!black, very thick] (0,-13.4) --
    node[above, font=\sffamily\tiny, text=red!60!black]
    {FINALIZE: consume UTxO \#1 + \#2 + \#3 atomically} (6,-13.4);
  \node[guidebox, fill=red!6, minimum width=4cm, minimum height=0.7cm,
    font=\sffamily\tiny] at (6,-14.2)
    {Single atomic transaction -- all or nothing};

  % Phase 6: Verify
  \draw[thick, purple!60] (-1.8,-15.0) -- (13.5,-15.0);
  \node[guidelabel, anchor=west, font=\sffamily\footnotesize\bfseries,
    text=purple!70!black] at (-1.8,-15.3) {Phase 6: On-Chain Verification};
  \node[guidebox, fill=gray!10, minimum width=1.8cm, minimum height=0.5cm,
    font=\sffamily\tiny] at (10,-15.8) {Any Verifier};
  \draw[guidearrow, purple!70] (10,-16.2) --
    node[above, font=\sffamily\tiny] {Query: who signed? when? valid credentials?}
    (1,-16.2);
  \draw[guidearrow, purple!70, dashed] (1,-16.8) --
    node[above, font=\sffamily\tiny] {Proof: PKHs, slots, doc hash, policy IDs}
    (10,-16.8);

  % Time arrow
  \draw[guidearrow, gray!50, very thick] (14,-0.5) -- (14,-17.5);
  \node[guidelabel, rotate=-90] at (14.5,-9) {Time};
\end{tikzpicture}
\caption{Complete six-phase multi-party signing sequence with three signers.
  Each signer creates an independent UTxO at the validator address (Phase~3),
  avoiding eUTxO concurrency conflicts.  The authority's finalize transaction
  (Phase~5) atomically consumes all deposits.}
\label{fig:six-phase-sequence}
\end{figure}


% ----------------------------------------------------------------------------
\section{The eUTxO Concurrency Solution}
\label{sec:eutxo-concurrency}

\subsection{Understanding the Problem}

Cardano uses the \textbf{extended UTxO (eUTxO) model}, which is
fundamentally different from Ethereum's account model.  In the account model,
a smart contract has a single ``account'' whose state is updated by
transactions.  Multiple transactions can modify the same account in the same
block (the EVM processes them sequentially within the block).  In the eUTxO
model, each UTxO is a discrete, one-time-use object: a transaction
\emph{consumes} one or more UTxOs and \emph{produces} new ones.  A UTxO can
be consumed exactly once---it is destroyed in the process.

This creates a concurrency challenge.  Suppose the system used a naive design
where all signers deposit into the \emph{same} UTxO at the validator address.
Signer~A sees UTxO~X (containing zero deposits) and builds a transaction that
consumes UTxO~X and produces UTxO~X$'$ (containing A's deposit).  Meanwhile,
Signer~B also sees UTxO~X and builds a transaction consuming it to produce
UTxO~X$''$ (containing B's deposit).  Both transactions reference UTxO~X as
an input.  But UTxO~X can only be consumed once.

When both transactions are submitted to the network, only one will succeed.
The other will fail with an error indicating that its input UTxO has already
been consumed.  This is the \textbf{eUTxO contention problem}: concurrent
transactions competing to spend the same UTxO cause all but one to fail.

In high-throughput DeFi applications on Cardano, this problem is well-known
and has motivated various batching and order-book designs.  For a credential
signing system, however, the problem is especially acute because signers are
\emph{expected} to act concurrently---telling three professionals ``please
sign one at a time and wait for on-chain confirmation between each'' is
operationally unacceptable.

\subsection{The Independent UTxO Solution}

The solution is elegant in its simplicity: \textbf{each signer creates their
own independent UTxO at the script address}.  There is no shared UTxO to
contend over.  Signer~A creates UTxO$_A$, Signer~B creates UTxO$_B$, and
Signer~C creates UTxO$_C$---all at the same validator address, but as
separate, non-conflicting UTxOs.  Since no two transactions reference the
same input, they can all succeed in the same block if submitted
simultaneously.

The authority then ``sweeps'' all per-signer UTxOs in a single finalization
transaction.  This finalize transaction has multiple inputs (one per signer
UTxO) and the authority's signature, and it is the only transaction that
needs to consume multiple UTxOs atomically.  Since only the authority submits
the finalize transaction, there is no concurrency conflict---only one party
is building this transaction.

\vspace{0.5cm}
\begin{figure}[H]
\centering
\begin{tikzpicture}[scale=0.82, every node/.style={transform shape}]
  % ---- LEFT SIDE: THE PROBLEM ----
  \node[font=\sffamily\bfseries, text=red!60!black] at (-4.5, 5.5)
    {THE PROBLEM: Shared UTxO};

  % Shared UTxO
  \node[guidebox, fill=red!6, minimum width=2.5cm, minimum height=1cm]
    (shared) at (-4.5, 4) {Shared UTxO at Validator};

  % Three signers trying to spend it
  \node[guidebox, fill=blue!8, minimum width=1.8cm] (sa) at (-7, 2) {Signer A};
  \node[guidebox, fill=green!8, minimum width=1.8cm] (sb) at (-4.5, 2) {Signer B};
  \node[guidebox, fill=orange!8, minimum width=1.8cm] (sc) at (-2, 2) {Signer C};

  \draw[guidearrow, red!60!black] (sa) -- (shared);
  \draw[guidearrow, red!60!black] (sb) -- (shared);
  \draw[guidearrow, red!60!black] (sc) -- (shared);

  % Conflict marker
  \node[font=\sffamily\scriptsize, text=red!70!black, text width=3cm,
    align=center] at (-4.5, 0.6)
    {Only ONE succeeds!\\Others fail with\\``UTxO already spent''};

  % X marks
  \node[font=\bfseries\Large, text=red!70] at (-7, 3.1) {$\times$};
  \node[font=\bfseries\Large, text=green!70!black] at (-4.5, 3.1) {\checkmark};
  \node[font=\bfseries\Large, text=red!70] at (-2, 3.1) {$\times$};

  % ---- RIGHT SIDE: THE SOLUTION ----
  \node[font=\sffamily\bfseries, text=green!60!black] at (5.5, 5.5)
    {THE SOLUTION: Independent UTxOs};

  % Three independent UTxOs
  \node[guidebox, fill=blue!8, minimum width=1.8cm] (sa2) at (3, 2)
    {Signer A};
  \node[guidebox, fill=green!8, minimum width=1.8cm] (sb2) at (5.5, 2)
    {Signer B};
  \node[guidebox, fill=orange!8, minimum width=1.8cm] (sc2) at (8, 2)
    {Signer C};

  % Independent UTxOs at validator
  \node[guidebox, fill=blue!5, minimum width=1.5cm, minimum height=0.7cm,
    font=\sffamily\tiny] (ua) at (3, 4) {UTxO$_A$};
  \node[guidebox, fill=green!5, minimum width=1.5cm, minimum height=0.7cm,
    font=\sffamily\tiny] (ub) at (5.5, 4) {UTxO$_B$};
  \node[guidebox, fill=orange!5, minimum width=1.5cm, minimum height=0.7cm,
    font=\sffamily\tiny] (uc) at (8, 4) {UTxO$_C$};

  \draw[guidearrow, green!60!black] (sa2) -- (ua);
  \draw[guidearrow, green!60!black] (sb2) -- (ub);
  \draw[guidearrow, green!60!black] (sc2) -- (uc);

  % All succeed
  \node[font=\bfseries\Large, text=green!70!black] at (3, 3.1) {\checkmark};
  \node[font=\bfseries\Large, text=green!70!black] at (5.5, 3.1) {\checkmark};
  \node[font=\bfseries\Large, text=green!70!black] at (8, 3.1) {\checkmark};

  % Authority sweep
  \node[guidebox, fill=purple!8, minimum width=3cm, minimum height=0.8cm]
    (auth) at (5.5, -0.5) {Authority -- FINALIZE};
  \draw[guidearrow, purple!70, thick] (ua.south) -- (auth.north west);
  \draw[guidearrow, purple!70, thick] (ub.south) -- (auth.north);
  \draw[guidearrow, purple!70, thick] (uc.south) -- (auth.north east);

  \node[font=\sffamily\scriptsize, text=green!60!black, text width=3.2cm,
    align=center] at (5.5, -1.6)
    {All signers succeed concurrently. Authority sweeps in one atomic tx.};

  % Validator address label
  \draw[dashed, gray!40, thick] (1.8, 3.4) rectangle (9.2, 4.7);
  \node[guidelabel, anchor=south] at (5.5, 4.7)
    {Validator Script Address};

  % Divider
  \draw[thick, gray!30] (0.5, -2) -- (0.5, 5.8);
\end{tikzpicture}
\caption{eUTxO concurrency: the shared-UTxO problem (left) versus the
  independent-UTxO solution (right).  When multiple signers target the same
  UTxO, only one transaction succeeds.  With independent UTxOs, all signers
  can deposit concurrently, and the authority sweeps all deposits in a single
  finalize transaction.}
\label{fig:eutxo-concurrency}
\end{figure}

\subsection{Why This Works in Cardano's Model}

The independent-UTxO pattern works because Cardano's eUTxO model treats each
UTxO as a separate object.  A validator script is identified by its script
hash, which determines the script address.  \emph{Any number of UTxOs} can
sit at the same script address, each with its own datum, its own tokens, and
its own spending conditions.  The validator script is invoked separately for
each UTxO being consumed, receiving the datum attached to \emph{that
specific UTxO} and the redeemer provided by the transaction.

This means the validator can enforce per-signer rules (``the datum's PKH
matches a required signer'') without needing to manage a shared state.
The ``global'' view---``have all required signers deposited?''---is checked
only in the finalize transaction, where the authority provides all UTxOs as
inputs and the script can verify completeness across the full input set.

This design pattern---per-user UTxOs at a shared script address, with a
privileged party performing the aggregation step---is a general solution
to eUTxO concurrency and is used in Cardano DEX order books, lending
protocols, and governance systems.  In the credential context, it maps
naturally to the real-world workflow: professionals act independently,
and an authority aggregates their actions into a final result.


% ----------------------------------------------------------------------------
\section{SignerDatum Structure}
\label{sec:signer-datum}

Each deposit UTxO at the validator address carries a \texttt{SignerDatum}---a
structured piece of data (a Plutus \texttt{PlutusData} type serialized as
CBOR) that creates an immutable record of the signing act.  The datum
contains five fields:

\begin{table}[H]
\centering
\caption{SignerDatum fields}
\label{tab:signer-datum}
\small
\begin{tabular}{@{}l l p{7cm}@{}}
\toprule
\textbf{Field} & \textbf{Type} & \textbf{Description} \\
\midrule
\texttt{signer\_pkh} & \texttt{PubKeyHash} (28 bytes) &
  The public key hash of the signer.  This is derived from the signer's
  Ed25519 verification key and serves as the on-chain identity binding.
  Combined with the transaction signature, it proves that the holder of
  the corresponding private key authorized this deposit. \\
\texttt{document\_hash} & \texttt{ByteString} (32 bytes) &
  The SHA-256 hash of the document being signed.  This binds the signature
  to a specific, unalterable document version.  Any party with the original
  document can recompute the hash and verify it matches. \\
\texttt{policy\_id} & \texttt{MintingPolicyHash} (28 bytes) &
  The policy ID under which the signer's credential tokens were minted.
  This allows verifiers to confirm which authority issued the signer's
  credentials, establishing provenance. \\
\texttt{validity\_expiry} & \texttt{POSIXTime} (slot number) &
  The expiry slot of the validity token deposited alongside this datum.
  Combined with \texttt{deposit\_slot}, this field proves that the signer's
  credentials were valid (not expired) at the time of signing. \\
\texttt{deposit\_slot} & \texttt{POSIXTime} (slot number) &
  The slot at which this deposit transaction was included in a block.
  This provides a precise timestamp of the signing act, anchored to
  Cardano's consensus-verified slot timeline. \\
\bottomrule
\end{tabular}
\end{table}

The combination of these five fields creates a self-contained proof of
credential-verified signing.  A verifier needs nothing beyond the Cardano
blockchain to confirm:  (1)~\texttt{signer\_pkh} identifies who signed;
(2)~\texttt{document\_hash} identifies what was signed;
(3)~\texttt{policy\_id} identifies which authority vouches for the signer;
(4)~\texttt{validity\_expiry} $>$ \texttt{deposit\_slot} proves the
credential was valid at signing time.

The datum is attached to the UTxO when it is created (Phase~3) and becomes
part of the blockchain's permanent record when the UTxO is consumed (Phase~5).
Because Cardano stores spent transaction data in the immutable ledger history,
the \texttt{SignerDatum} is preserved indefinitely---even after the UTxO
itself has been spent.  This is the mechanism by which the system provides
permanent, tamper-proof signing records.


% ============================================================================
% PART B: PRIVACY AND COMPLIANCE
% ============================================================================

\vspace{0.5cm}
\noindent\rule{\textwidth}{1.2pt}
\begin{center}
  {\Large\sffamily\bfseries Part B: Privacy and Compliance}
\end{center}
\noindent\rule{\textwidth}{1.2pt}
\vspace{0.3cm}


% ----------------------------------------------------------------------------
\section{The Privacy Architecture}
\label{sec:privacy-architecture}

A blockchain is, by design, a public, transparent, and immutable ledger.
Every transaction, every token transfer, every datum is visible to anyone
running a node or querying an explorer.  This transparency is precisely
what makes blockchain valuable for credential verification---but it creates
an immediate tension with data privacy requirements.  You cannot put a
person's name, license number, home address, and date of birth on a public
blockchain and simultaneously comply with data protection laws.

The system resolves this tension through a strict \textbf{architectural
separation} between on-chain and off-chain data.  The principle is simple
but rigorously applied: \emph{anything that touches the blockchain must be
incapable of identifying a natural person without access to a separate,
private mapping.}

\subsection{On-Chain Data (Public, Immutable)}

The following data elements exist on the Cardano blockchain and are visible
to any observer:

\begin{itemize}[leftmargin=*, nosep, itemsep=4pt]
  \item \textbf{Policy IDs.}  These are Blake2b-224 hashes of minting
    scripts.  A policy ID like \texttt{a1b2c3...f4e5} reveals nothing about
    the authority that created it unless the observer already knows the
    mapping between policy IDs and authorities.
  \item \textbf{Token asset names.}  Names like \texttt{LIC\_PE\_001} or
    \texttt{SIG\_TOKEN} are visible on-chain but contain no personal data.
    The naming convention uses role-based prefixes, not personal identifiers.
  \item \textbf{Transaction signatures (PKHs).}  Public key hashes are
    28-byte values derived from Ed25519 verification keys.  A PKH like
    \texttt{8a3f...c7d2} is a pseudonym---it identifies a \emph{key pair},
    not a \emph{person}.  Without the off-chain mapping, an observer cannot
    determine whose key it is.
  \item \textbf{CIP-25/68 metadata fields.}  The on-chain metadata attached
    to license NFTs includes credential type, jurisdiction code, issue date,
    status, and policy-specific fields.  These are designed to be
    \emph{functionally useful} (a verifier can confirm ``this is a PE
    license issued in California that is currently active'') without being
    \emph{personally identifying} (the metadata does not include the
    licensee's name or license number).
  \item \textbf{Slot timestamps.}  Cardano's slot numbers provide temporal
    ordering and can be converted to UTC timestamps.  These record
    \emph{when} events occurred but not \emph{who} was involved (beyond
    the pseudonymous PKH).
\end{itemize}

\subsection{Off-Chain Data (Private, Erasable)}

Personal data that connects blockchain pseudonyms to real-world identities
is stored exclusively in an off-chain database:

\begin{itemize}[leftmargin=*, nosep, itemsep=4pt]
  \item \textbf{Full names} of licensees, signers, and authority contacts
  \item \textbf{License numbers} as assigned by the issuing jurisdiction
  \item \textbf{Physical addresses} and contact information
  \item \textbf{Email addresses} and phone numbers
  \item \textbf{PKH-to-identity mappings} that link on-chain pseudonyms
    to real people
\end{itemize}

This data is stored in an \textbf{AES-256-GCM encrypted SQLite database}.
The encryption ensures that even if an attacker gains access to the database
file, the contents are unreadable without the encryption key.

\vspace{0.5cm}
\begin{figure}[H]
\centering
\begin{tikzpicture}[scale=0.82, every node/.style={transform shape}]
  % On-chain section
  \draw[thick, blue!40, rounded corners=8pt, fill=blue!3]
    (-5.5, 0) rectangle (5.5, 6.5);
  \node[font=\sffamily\bfseries, text=blue!70!black] at (0, 6.1)
    {ON-CHAIN (Public, Immutable)};
  \node[font=\sffamily\small, text=blue!60!black] at (0, 5.6)
    {Cardano Blockchain --- visible to all observers};

  \node[guidebox, fill=blue!8, minimum width=2.2cm, minimum height=0.8cm,
    font=\sffamily\scriptsize, text width=2cm] (pid) at (-3.2, 4.5) {Policy IDs (Blake2b hashes)};
  \node[guidebox, fill=blue!8, minimum width=2.2cm, minimum height=0.8cm,
    font=\sffamily\scriptsize, text width=2cm] (asset) at (0, 4.5) {Asset Names (role prefixes)};
  \node[guidebox, fill=blue!8, minimum width=2.2cm, minimum height=0.8cm,
    font=\sffamily\scriptsize, text width=2cm] (txsig) at (3.2, 4.5) {Tx Signatures (PKH pseudonyms)};

  \node[guidebox, fill=blue!8, minimum width=2.2cm, minimum height=0.8cm,
    font=\sffamily\scriptsize] (meta) at (-1.6, 3) {CIP-25/68 Metadata};
  \node[guidebox, fill=blue!8, minimum width=2.2cm, minimum height=0.8cm,
    font=\sffamily\scriptsize] (slots) at (1.6, 3) {Slot Timestamps};

  \node[guidebox, fill=yellow!8, minimum width=4cm, minimum height=0.7cm,
    font=\sffamily\scriptsize\itshape] at (0, 1.5)
    {No names, no license numbers, no addresses, no contact info};

  % Separator
  \draw[very thick, red!50!black, dashed] (-6, -0.5) -- (6, -0.5);
  \node[font=\sffamily\bfseries, text=red!60!black,
    fill=white, inner sep=3pt] at (0, -0.5)
    {PRIVACY BOUNDARY};

  % Off-chain section
  \draw[thick, green!40, rounded corners=8pt, fill=green!3]
    (-5.5, -6.5) rectangle (5.5, -1.2);
  \node[font=\sffamily\bfseries, text=green!70!black] at (0, -1.6)
    {OFF-CHAIN (Private, Erasable)};
  \node[font=\sffamily\small, text=green!60!black] at (0, -2.1)
    {AES-256-GCM Encrypted SQLite Database};

  \node[guidebox, fill=green!8, minimum width=2.2cm, minimum height=0.8cm,
    font=\sffamily\scriptsize] (names) at (-3.2, -3.2) {Full Names};
  \node[guidebox, fill=green!8, minimum width=2.2cm, minimum height=0.8cm,
    font=\sffamily\scriptsize] (licnum) at (0, -3.2) {License Numbers};
  \node[guidebox, fill=green!8, minimum width=2.2cm, minimum height=0.8cm,
    font=\sffamily\scriptsize] (addr) at (3.2, -3.2) {Addresses / Contact Info};

  \node[guidebox, fill=green!8, minimum width=2.2cm, minimum height=0.8cm,
    font=\sffamily\scriptsize] (email) at (-1.6, -4.6) {Email \& Phone};
  \node[guidebox, fill=green!8, minimum width=2.2cm, minimum height=0.8cm,
    font=\sffamily\scriptsize] (mapping) at (1.6, -4.6) {PKH $\leftrightarrow$ Identity Map};

  % Lock icon area
  \node[guidebox, fill=red!6, minimum width=4cm, minimum height=0.7cm,
    font=\sffamily\scriptsize\bfseries] at (0, -5.8)
    {Encryption key controls access --- Destroy key = erase data (GDPR Art.\ 17)};

  % Connection arrows
  \draw[guidearrow, dashed, gray!50] (txsig.south) -- ++(0,-2.5)
    node[midway, right, font=\sffamily\tiny, text=gray]
    {PKH links to identity only via encrypted mapping};
\end{tikzpicture}
\caption{Privacy architecture: strict separation between public on-chain data
  (policy IDs, asset names, PKHs, metadata, timestamps) and private off-chain
  data (names, license numbers, addresses, contact information).  The PKH-to-identity
  mapping exists only in the encrypted off-chain database.}
\label{fig:privacy-split}
\end{figure}


% ----------------------------------------------------------------------------
\section{AES-256-GCM Encryption}
\label{sec:aes256gcm}

The off-chain database and the authority's wallet keys at rest are protected
by \textbf{AES-256-GCM} encryption.  Understanding what this means requires
unpacking the acronym:

\textbf{AES} stands for \emph{Advanced Encryption Standard}, a symmetric-key
encryption algorithm adopted by NIST in 2001 (FIPS~197).  ``Symmetric-key''
means the same key is used for both encryption and decryption.  AES operates
on 128-bit blocks of data, transforming plaintext into ciphertext through a
series of substitution, permutation, and mixing operations.

\textbf{256} refers to the key length in bits.  A 256-bit key provides
$2^{256}$ possible key values---a number so large ($\approx 1.16 \times
10^{77}$) that brute-force enumeration is considered computationally
infeasible for the foreseeable future, including against quantum computers
using Grover's algorithm (which reduces the effective security to 128 bits,
still considered sufficient).

\textbf{GCM} stands for \emph{Galois/Counter Mode}, an authenticated
encryption mode that provides two guarantees simultaneously:

\begin{enumerate}[nosep]
  \item \textbf{Confidentiality.}  The ciphertext reveals nothing about the
    plaintext to anyone who does not possess the key.  GCM achieves this
    using CTR (Counter) mode, which turns AES into a stream cipher by
    encrypting sequential counter values and XOR-ing the result with the
    plaintext.
  \item \textbf{Integrity and Authentication.}  GCM produces an
    authentication tag (typically 128 bits) alongside the ciphertext.  This
    tag is computed using Galois field multiplication over the ciphertext
    blocks.  When decrypting, the recipient recomputes the tag and compares
    it with the stored tag.  If even a single bit of the ciphertext has been
    tampered with, the tags will not match, and the decryption is rejected.
\end{enumerate}

The combination of confidentiality and integrity is critical for the
credential system.  Without integrity protection, an attacker who gains
access to the encrypted database could modify ciphertext bytes (even
without knowing the key) and potentially cause the decrypted data to
contain corrupted-but-accepted values.  GCM prevents this: any modification
to the ciphertext is detected and rejected.

In the credential system, AES-256-GCM is used in two contexts:

\begin{itemize}[nosep]
  \item \textbf{Authority wallet keys at rest.}  The authority's Ed25519
    signing key is encrypted before being stored in the SQLite database.
    The encryption key is derived from a passphrase or stored in a separate
    key management system.  This ensures that database compromise alone
    does not expose the authority's signing capability.
  \item \textbf{Off-chain PII database.}  Personal data (names, license
    numbers, addresses) is encrypted at the field level or database level.
    Per-credential encryption keys enable granular erasure: destroying one
    key renders one person's data permanently inaccessible without affecting
    other records.
\end{itemize}


% ----------------------------------------------------------------------------
\section{GDPR and Crypto-Shredding}
\label{sec:gdpr-crypto-shredding}

\subsection{The GDPR Dilemma}

The General Data Protection Regulation (GDPR), enforced across the European
Union since May 2018, establishes a comprehensive framework for personal data
protection.  Article~17 is the provision most directly in tension with
blockchain technology:

\begin{quote}
\textit{``The data subject shall have the right to obtain from the controller
the erasure of personal data concerning him or her without undue delay, and
the controller shall have the obligation to erase personal data without
undue delay\ldots''}  --- GDPR Article 17(1)
\end{quote}

This is the \textbf{right to erasure}, commonly called the ``right to be
forgotten.''  It requires that upon request, a data controller must delete
all personal data associated with the requesting individual.

Blockchain, however, is \emph{immutable by design}.  Once a transaction is
confirmed and finalized, it cannot be altered or deleted.  This is not a
bug---it is the core value proposition.  Credential verification depends on
the guarantee that a licensing board's attestation, once published, cannot
be retroactively modified or erased.

The tension is stark: GDPR says ``delete it,'' and blockchain says ``it
cannot be deleted.''  The European Data Protection Board (EDPB) has made
clear in its April 2025 guidelines that technical impossibility is not an
acceptable excuse for non-compliance, with penalties reaching up to 20
million EUR or 4\% of global annual turnover, whichever is higher.

\subsection{The Solution: Crypto-Shredding}

The credential system resolves this dilemma through a technique called
\textbf{crypto-shredding} (also known as cryptographic erasure).  The
approach is architecturally simple but legally significant:

\begin{enumerate}[nosep]
  \item \textbf{Personal data is never stored on-chain.}  As described in
    Section~\ref{sec:privacy-architecture}, names, license numbers,
    addresses, and other personally identifiable information (PII) exist
    only in the encrypted off-chain database.
  \item \textbf{Each credential's PII is encrypted with a per-credential
    key.}  Rather than encrypting the entire database with a single key,
    the system uses individual encryption keys for each person's data.
    This enables selective erasure without affecting other records.
  \item \textbf{To ``erase'' data, destroy the encryption key.}  When a
    data subject exercises their Article~17 right, the system destroys the
    encryption key associated with their data.  The encrypted ciphertext
    remains in the database, but without the key, it is computationally
    indistinguishable from random noise.  The data is effectively and
    permanently erased.
  \item \textbf{On-chain hashes become meaningless.}  The on-chain data
    consists of cryptographic hashes (policy IDs, PKHs, asset names) and
    metadata that does not contain PII.  A PKH like \texttt{8a3f...c7d2}
    cannot be reversed to recover the original public key (hash functions
    are one-way), and even if it could, the public key does not reveal
    the person's identity without the destroyed off-chain mapping.
\end{enumerate}

This approach aligns with guidance from two regulatory bodies:

\begin{itemize}[nosep]
  \item \textbf{French CNIL (2018).}  The Commission Nationale de
    l'Informatique et des Libert\'{e}s published guidelines recognizing
    that storing only hashed or encrypted data on-chain, combined with
    off-chain key management and destruction capabilities, constitutes
    a GDPR-compliant architecture.
  \item \textbf{EDPB Framework (2025).}  The European Data Protection Board's
    updated guidelines acknowledge crypto-shredding as a valid technical
    measure for blockchain systems, provided that the on-chain data cannot
    be used to re-identify individuals after key destruction.
\end{itemize}

\vspace{0.5cm}
\begin{figure}[H]
\centering
\begin{tikzpicture}[scale=0.82, every node/.style={transform shape}]
  % Title
  \node[font=\sffamily\bfseries\large, text=purple!70!black] at (0, 8.5)
    {Crypto-Shredding Process};

  % ---- BEFORE: Normal Operation ----
  \node[font=\sffamily\bfseries, text=blue!70!black] at (-5, 7.5)
    {BEFORE: Normal Operation};

  % On-chain
  \node[guidebox, fill=blue!8, minimum width=3.5cm, minimum height=1.2cm,
    font=\sffamily\scriptsize, text width=3cm] (onchain1) at (-5, 5.8)
    {On-Chain\\[2pt]PKH: \texttt{8a3f...c7d2}\\Policy: \texttt{a1b2...f4e5}\\Status: active};

  % Mapping arrow
  \draw[guidearrow, thick, blue!60] (-5, 5) -- (-5, 3.8);
  \node[guidebox, fill=yellow!8, minimum width=3cm, minimum height=0.7cm,
    font=\sffamily\scriptsize\bfseries] (key1) at (-5, 3.4)
    {Encryption Key $K_i$};

  % Off-chain
  \node[guidebox, fill=green!8, minimum width=3.5cm, minimum height=1.2cm,
    font=\sffamily\scriptsize, text width=3cm] (offchain1) at (-5, 2)
    {Off-Chain (Encrypted)\\[2pt]Name: Jane Smith, PE\\License: PE-2026-042\\Address: 123 Main St};

  \draw[guidearrow, thick, green!60] (key1) -- (offchain1.north);

  % ---- AFTER: Crypto-Shredding ----
  \node[font=\sffamily\bfseries, text=red!70!black] at (5, 7.5)
    {AFTER: Key Destruction};

  % On-chain (unchanged)
  \node[guidebox, fill=blue!8, minimum width=3.5cm, minimum height=1.2cm,
    font=\sffamily\scriptsize, text width=3cm] (onchain2) at (5, 5.8)
    {On-Chain (unchanged)\\[2pt]PKH: \texttt{8a3f...c7d2}\\Policy: \texttt{a1b2...f4e5}\\Status: active};

  % Key destroyed
  \draw[guidearrow, thick, red!60, dashed] (5, 5) -- (5, 3.8);
  \node[guidebox, fill=red!12, minimum width=3cm, minimum height=0.7cm,
    font=\sffamily\scriptsize\bfseries, text=red!70!black] (key2) at (5, 3.4)
    {Key $K_i$ DESTROYED};

  % Draw X over key
  \draw[very thick, red!70] (3.8, 3.1) -- (6.2, 3.7);
  \draw[very thick, red!70] (3.8, 3.7) -- (6.2, 3.1);

  % Off-chain (unreadable)
  \node[guidebox, fill=gray!15, minimum width=3.5cm, minimum height=1.2cm,
    font=\sffamily\scriptsize, text width=3cm, text=gray!50] (offchain2) at (5, 2)
    {Off-Chain (Ciphertext)\\[2pt]\texttt{7f2a 9bc1 4e3d ...}\\Indistinguishable\\from random noise};

  \draw[guidearrow, thick, gray!40, dashed] (key2) -- (offchain2.north);

  % Arrow from before to after
  \draw[-{Latex[length=4mm]}, very thick, red!50!black]
    (-1.5, 4) -- node[above, font=\sffamily\small\bfseries, text=red!60!black]
    {GDPR Art.\ 17 Erasure Request} (1.5, 4);

  % Result annotation
  \node[guidebox, fill=purple!8, minimum width=8cm, minimum height=1cm,
    font=\sffamily\scriptsize, text width=7.5cm] at (0, 0)
    {\textbf{Result:} On-chain hashes are meaningless without the mapping.
     Hash functions are one-way (cannot reverse PKH to identity).
     Encrypted off-chain data is permanently inaccessible.
     GDPR compliance achieved without modifying the blockchain.};
\end{tikzpicture}
\caption{Crypto-shredding process.  Left: during normal operation, the
  encryption key $K_i$ connects on-chain pseudonyms to off-chain personal
  data.  Right: after a GDPR erasure request, destroying $K_i$ renders the
  off-chain data permanently unreadable while on-chain hashes remain but
  become unlinkable to any natural person.}
\label{fig:crypto-shredding}
\end{figure}

\subsection{Why Hashes Cannot Be Reversed}

A common concern is: ``If the PKH is derived from a public key, can't
someone reverse the hash to recover the public key, and then use the public
key to identify the person?''  The answer is no, for two reasons:

First, Blake2b-224 (the hash function used for Cardano PKHs) is a
cryptographic hash function with preimage resistance.  Given a hash output
$h$, finding any input $x$ such that $\text{Blake2b-224}(x) = h$ requires
approximately $2^{224}$ operations---a number vastly exceeding the
computational capacity of all computers ever built.

Second, even if an attacker could somehow recover the public key from the
PKH, the public key is a 32-byte Ed25519 verification key.  This key does
not inherently contain any personal information.  The mapping from public
key to person exists only in the off-chain database---which has been
crypto-shredded.

The combination of one-way hash functions and off-chain key destruction
creates what privacy engineers call \textbf{unlinkability}: the on-chain
data points cannot be linked to a natural person after crypto-shredding,
satisfying the GDPR's erasure requirement in substance.


% ============================================================================
% PART C: LEGAL FRAMEWORK
% ============================================================================

\vspace{0.5cm}
\noindent\rule{\textwidth}{1.2pt}
\begin{center}
  {\Large\sffamily\bfseries Part C: Legal Framework}
\end{center}
\noindent\rule{\textwidth}{1.2pt}
\vspace{0.3cm}


% ----------------------------------------------------------------------------
\section{ESIGN Act (US Federal)}
\label{sec:esign}

The Electronic Signatures in Global and National Commerce Act (15~U.S.C.
\S\S~7001--7031), enacted in 2000, is the foundational US federal law
governing electronic signatures.  Its central provision is deceptively
simple: \emph{an electronic signature cannot be denied legal effect solely
because it is in electronic form}.  The Act is deliberately
technology-neutral---it does not prescribe \emph{how} an electronic
signature must be implemented, only that it must satisfy four requirements.

\subsection{Requirement 1: Intent to Sign}

The signer must demonstrate a deliberate intent to sign the document.  In
traditional e-signature tools, this is typically satisfied by clicking an
``I agree'' button.  In the credential system, intent is demonstrated by
the deliberate invocation of the \texttt{sign\_document()} function, which
requires the signer to:

\begin{itemize}[nosep]
  \item Identify the specific document (by hash) they are signing
  \item Decrypt their private signing key (requiring passphrase entry
    or hardware wallet interaction)
  \item Construct and submit an on-chain transaction that deposits tokens
    to the validator address
\end{itemize}

This multi-step process---far more deliberate than clicking a button---provides
strong evidence of intent.  A signer cannot accidentally deposit tokens to
a validator address; the process requires explicit, informed action at
multiple stages.

\subsection{Requirement 2: Consent to Conduct Business Electronically}

The signer must have consented to conducting business electronically.  The
credential system satisfies this through the onboarding process: when a
professional registers for the system and receives their wallet, credential
tokens, and access credentials, they explicitly accept terms of service that
include consent to electronic transactions.  This consent is recorded in the
off-chain database with a timestamp and is presented to the signer before
any wallet generation or token minting occurs.

\subsection{Requirement 3: Association of Signature with Record}

The electronic signature must be logically associated with the record being
signed.  The credential system achieves this through cryptographic binding
at three levels:

\begin{enumerate}[nosep]
  \item \textbf{SHA-256 hash binding.}  The \texttt{SignerDatum} includes
    the SHA-256 hash of the document, creating a mathematical link between
    the signature and the exact document version.
  \item \textbf{Transaction binding.}  The deposit transaction is signed
    with the signer's Ed25519 private key, cryptographically proving that
    the holder of that key authorized this specific transaction.
  \item \textbf{UTxO binding.}  The deposited tokens (signature token +
    validity token) are locked at the validator address with the
    \texttt{SignerDatum}, creating an inseparable on-chain association
    between the signer's identity, the document, and the signing act.
\end{enumerate}

This triple binding is significantly stronger than the ``email address
clicked a link'' association provided by most commercial e-signature tools.

\subsection{Requirement 4: Record Retention}

The signed record must be retained and reproducible.  The credential system
satisfies this through blockchain immutability: once the finalize transaction
is confirmed, the complete signing record---who signed, what they signed,
when they signed, and what credentials they held---is stored across
thousands of Cardano nodes worldwide.  The data cannot be altered, deleted,
or lost (short of a catastrophic failure of the entire Cardano network).
This level of retention surpasses any centralized storage system in both
durability and tamper-resistance.


% ----------------------------------------------------------------------------
\section{UETA (US State Law)}
\label{sec:ueta}

The Uniform Electronic Transactions Act (UETA) was drafted by the Uniform
Law Commission in 1999 and has been adopted by \textbf{49 states plus the
District of Columbia} (New York is the sole non-adopter, though it has its
own Electronic Signatures and Records Act).  UETA provides the state-level
legal foundation for electronic signatures, complementing the federal ESIGN
Act.

UETA's key contribution is Section~9, which establishes that an electronic
signature is ``attributable to a person'' if it can be shown by a
``preponderance of the evidence'' that the person performed the act.
Cardano's Ed25519 cryptography provides exceptionally strong attribution:
possessing the private key corresponding to a PKH is, in practice,
conclusive evidence of authorship.  The probability of forging an Ed25519
signature without the private key is approximately $2^{-128}$---effectively
zero.

\subsection{Blockchain-Specific State Amendments}

Nine states have gone beyond basic UETA adoption to enact
\textbf{blockchain-specific amendments} that explicitly recognize blockchain
signatures and smart contracts.  Two are particularly significant for the
credential system:

\textbf{Arizona HB 2417 (2017).}  Arizona was the first state to explicitly
recognize blockchain signatures.  A.R.S.~\S44-7061 provides that
``a signature that is secured through blockchain technology is considered
to be in an electronic form and to be an electronic signature'' and that
``a record or contract that is secured through blockchain technology is
considered to be in an electronic form and to be an electronic record.''
This statute directly validates the credential system's on-chain signing
mechanism: token deposits at a validator address, secured by Ed25519
cryptography on the Cardano blockchain, are electronic signatures under
Arizona law.

\textbf{Vermont 12 V.S.A. \S1913 (2016).}  Vermont's statute creates a
\textbf{rebuttable presumption of authenticity} for blockchain records.
Specifically, a ``digital record electronically registered in a blockchain''
is ``self-authenticating'' under the Vermont Rules of Evidence, provided that
the record is accompanied by a ``written declaration of a qualified person''
stating that the record is accurate.  This is an evidentiary advantage:
in court, the party introducing a blockchain-verified credential does not
bear the initial burden of proving its authenticity---the opposing party must
affirmatively rebut the presumption.  For the credential system, this means
that a finalized work product with on-chain signing records carries
automatic evidentiary weight in Vermont courts.

The remaining seven states with blockchain-specific provisions---Nevada,
Tennessee, South Dakota, Wyoming, Illinois, Colorado, and others---provide
varying degrees of explicit recognition for blockchain records and smart
contracts.  While the specific language differs, the trend is clear:
US state law is moving toward explicit statutory recognition of blockchain
as a valid medium for legally binding records and signatures.


% ----------------------------------------------------------------------------
\section{eIDAS 2.0 (European Union)}
\label{sec:eidas}

The European Union's electronic identification and trust services regulation,
known as eIDAS, is the most comprehensive legal framework for electronic
signatures globally.  The revised regulation, eIDAS~2.0 (Regulation (EU)
2024/1183), entered into force in 2024 and introduces provisions
specifically relevant to blockchain-based systems.

\subsection{Three Signature Tiers}

eIDAS~2.0 defines three tiers of electronic signatures, each with increasing
legal weight and technical requirements:

\textbf{Tier 1: Simple Electronic Signature (SES).}  The broadest category,
encompassing any ``data in electronic form which is attached to or logically
associated with other data in electronic form and which is used by the
signatory to sign.''  A typed name at the bottom of an email qualifies.
There are no technical requirements beyond basic logical association.  SES
signatures are admissible as evidence but carry no presumption of validity.

\textbf{Tier 2: Advanced Electronic Signature (AES).}  An AES must satisfy
four requirements defined in Article~26:

\begin{enumerate}[nosep]
  \item \textbf{Uniquely linked to the signatory.}  Each signer's Ed25519
    key pair is unique---no two signers share the same key.  The PKH
    derived from the verification key is a unique 28-byte identifier.
  \item \textbf{Capable of identifying the signatory.}  The PKH-to-identity
    mapping in the off-chain database links the cryptographic identity to
    a natural person.  The CIP-68 Reference Token's metadata provides
    additional identification context.
  \item \textbf{Created using data under the signatory's sole control.}
    The Ed25519 private key, stored encrypted in the signer's wallet
    (or in a hardware wallet via CIP-30), is under the signer's exclusive
    control.  No other party---including the authority---possesses the
    signer's private key.
  \item \textbf{Linked to the data signed in such a way that any
    subsequent change is detectable.}  The SHA-256 document hash in the
    \texttt{SignerDatum}, combined with the transaction signature, ensures
    that any modification to the signed document would produce a different
    hash, immediately detectable by any verifier.
\end{enumerate}

The credential system natively satisfies \textbf{all four AES requirements}
without requiring any modifications or additional infrastructure.  This is
a significant legal result: AES signatures carry a presumption of validity
across all EU member states and are admissible as evidence with enhanced
legal weight.

\textbf{Tier 3: Qualified Electronic Signature (QES).}  A QES is an AES
that is additionally:
\begin{itemize}[nosep]
  \item Created by a \textbf{Qualified Electronic Signature Creation
    Device (QSCD)}---a certified hardware device that protects the
    signing key
  \item Based on a \textbf{Qualified Certificate} issued by a
    \textbf{Qualified Trust Service Provider (QTSP)} that has been
    audited and approved by an EU member state's supervisory body
\end{itemize}

QES signatures have the \emph{legal equivalent of handwritten signatures}
across all EU member states.  The credential system does not natively
achieve QES status because it does not use a certified QSCD or a
QTSP-issued certificate.  However, a QTSP could integrate the credential
system's on-chain verification with its own qualified certificate
infrastructure to provide QES-level signatures.

\subsection{Qualified Electronic Ledger (Article 45l)}

eIDAS~2.0 introduces a new trust service category: the \textbf{Qualified
Electronic Ledger}.  Article~45l defines this as a distributed ledger
operated by a QTSP that provides ``proof of the existence and integrity of
electronic data recorded therein.''  This provision was specifically designed
to accommodate blockchain technology within the EU trust services framework.

For the credential system, Article~45l provides a clear pathway: if the
Cardano blockchain (or a specific Cardano-based service) obtains QTSP
status for ledger operation, then credentials recorded on it would carry
the full legal weight of a qualified trust service.  Cardano's formally
verified Ouroboros consensus protocol, with its published security proofs
in peer-reviewed venues, provides a strong technical foundation for such
qualification.

\vspace{0.5cm}
\begin{figure}[H]
\centering
\begin{tikzpicture}[scale=0.82, every node/.style={transform shape}]
  % Title
  \node[font=\sffamily\bfseries\large, text=purple!70!black] at (0, 9)
    {eIDAS 2.0 Electronic Signature Tiers};

  % Tier 1: SES
  \draw[thick, gray!40, rounded corners=6pt, fill=gray!5]
    (-6, 5.5) rectangle (6, 8.2);
  \node[font=\sffamily\bfseries, text=gray!60!black, anchor=west]
    at (-5.5, 7.8) {Tier 1: Simple Electronic Signature (SES)};
  \node[font=\sffamily\small, text=gray!50!black, text width=10cm,
    anchor=west] at (-5.5, 7.2)
    {Any data in electronic form logically associated with signed data.};
  \node[guidebox, fill=gray!10, minimum width=3cm, font=\sffamily\scriptsize]
    at (-3.5, 6.2) {Typed name in email};
  \node[guidebox, fill=gray!10, minimum width=3cm, font=\sffamily\scriptsize]
    at (0.5, 6.2) {Click-to-sign button};
  \node[guidebox, fill=gray!10, minimum width=3cm, font=\sffamily\scriptsize]
    at (4.5, 6.2) {Scanned signature image};
  \node[font=\sffamily\scriptsize\itshape, text=gray!50!black, anchor=west]
    at (-5.5, 5.7) {Legal weight: Admissible, but no presumption of validity.};

  % Tier 2: AES
  \draw[thick, blue!50, rounded corners=6pt, fill=blue!3]
    (-6, 1.2) rectangle (6, 5.2);
  \node[font=\sffamily\bfseries, text=blue!70!black, anchor=west]
    at (-5.5, 4.8) {Tier 2: Advanced Electronic Signature (AES)};

  \node[guidebox, fill=blue!8, minimum width=4.8cm, minimum height=0.6cm,
    font=\sffamily\scriptsize] at (-3, 3.9)
    {1. Uniquely linked to signatory (Ed25519 key)};
  \node[guidebox, fill=blue!8, minimum width=4.8cm, minimum height=0.6cm,
    font=\sffamily\scriptsize] at (-3, 3.1)
    {2. Capable of identifying signatory (PKH map)};
  \node[guidebox, fill=blue!8, minimum width=4.8cm, minimum height=0.6cm,
    font=\sffamily\scriptsize] at (-3, 2.3)
    {3. Sole control of signing data (private key)};
  \node[guidebox, fill=blue!8, minimum width=4.8cm, minimum height=0.6cm,
    font=\sffamily\scriptsize] at (-3, 1.5)
    {4. Change detection (SHA-256 hash binding)};

  \node[guidebox, fill=green!12, minimum width=3.8cm, minimum height=2cm,
    font=\sffamily\scriptsize\bfseries, text=green!50!black] at (3.5, 2.7)
    {This credential system meets ALL 4 AES requirements natively};

  % Tier 3: QES
  \draw[thick, purple!50, rounded corners=6pt, fill=purple!3]
    (-6, -2.8) rectangle (6, 0.9);
  \node[font=\sffamily\bfseries, text=purple!70!black, anchor=west]
    at (-5.5, 0.5) {Tier 3: Qualified Electronic Signature (QES)};
  \node[font=\sffamily\small, text=purple!50!black, text width=10cm,
    anchor=west] at (-5.5, -0.1)
    {AES + Qualified Certificate + Qualified Signature Creation Device};

  \node[guidebox, fill=purple!8, minimum width=3.2cm, font=\sffamily\scriptsize]
    at (-3.5, -1) {QSCD (certified hardware)};
  \node[guidebox, fill=purple!8, minimum width=3.2cm, font=\sffamily\scriptsize]
    at (0.5, -1) {QTSP-issued certificate};
  \node[guidebox, fill=purple!8, minimum width=3.2cm, font=\sffamily\scriptsize]
    at (4.5, -1) {EU supervisory body audit};

  \node[font=\sffamily\scriptsize\itshape, text=purple!50!black, anchor=west]
    at (-5.5, -2)
    {Legal weight: Equivalent to handwritten signature across all EU states.};
  \node[font=\sffamily\scriptsize\bfseries, text=orange!60!black, anchor=west]
    at (-5.5, -2.5)
    {Pathway: QTSP integration + Article 45l Qualified Electronic Ledger};

  % Arrow showing progression
  \draw[guidearrow, very thick, gray!50] (-7, 6.8) -- (-7, 3.2)
    -- (-7, -1);
  \node[guidelabel, rotate=90, anchor=south] at (-7.4, 3)
    {Increasing legal weight $\longrightarrow$};
\end{tikzpicture}
\caption{eIDAS~2.0 electronic signature tiers.  The credential system
  natively meets all four Advanced Electronic Signature (AES) requirements.
  Qualified Electronic Signature (QES) status requires additional QTSP
  certification and hardware security modules, providing a clear upgrade
  pathway via Article~45l.}
\label{fig:eidas-tiers}
\end{figure}


% ----------------------------------------------------------------------------
\section{Putting It All Together: Compliance Summary}
\label{sec:compliance-summary}

The multi-party signing workflow, privacy architecture, and legal framework
are not independent concerns---they are three facets of a single design
problem.  The signing workflow must produce legally admissible records.
The privacy architecture must not undermine the evidentiary value of those
records.  The legal framework must be satisfied without compromising the
technical integrity of the system.

Table~\ref{tab:compliance-matrix} summarizes how each component of the system
maps to specific legal requirements across all three major frameworks.

\begin{table}[H]
\centering
\caption{Compliance matrix: system components mapped to legal requirements}
\label{tab:compliance-matrix}
\small
\begin{tabular}{@{}p{3cm} p{3.5cm} p{3.5cm} p{3.5cm}@{}}
\toprule
\textbf{System Component} & \textbf{ESIGN Act} & \textbf{UETA / State Law} & \textbf{eIDAS 2.0} \\
\midrule
Ed25519 key pair &
  Association (Req.~3) &
  Attribution (Sec.~9) &
  AES Reqs.~1, 3 \\
\texttt{sign\_document()} invocation &
  Intent (Req.~1) &
  Intent to sign &
  Signatory action \\
Onboarding consent &
  Consent (Req.~2) &
  Agreement to e-records &
  N/A \\
SHA-256 hash in datum &
  Association (Req.~3) &
  Record integrity &
  AES Req.~4 \\
Blockchain immutability &
  Retention (Req.~4) &
  Record preservation &
  Art.~45l ledger \\
CIP-68 credential status &
  N/A &
  Real-time verification &
  Trust service \\
Crypto-shredding &
  N/A &
  N/A &
  GDPR Art.~17 \\
PKH-to-identity mapping &
  N/A &
  N/A &
  AES Req.~2 \\
\bottomrule
\end{tabular}
\end{table}

The credential system achieves a rare balance: it provides stronger
cryptographic guarantees than traditional e-signature tools (SHA-256 hash
binding, Ed25519 signatures, immutable on-chain records), stronger credential
verification than any existing signing platform (real-time on-chain license
status), and GDPR compliance through architectural separation and
crypto-shredding---all while satisfying the requirements of three major
legal frameworks governing electronic signatures.

This compliance is not incidental; it is the result of architectural
decisions made at the system's foundation.  The separation of on-chain and
off-chain data, the use of PKHs as pseudonyms, the per-credential
encryption keys, and the independent-UTxO signing pattern were all
designed with legal requirements as first-class constraints alongside
technical requirements.

\section*{Chapter Summary}

\begin{itemize}[leftmargin=*, itemsep=4pt]
  \item Multi-party document signing uses a six-phase workflow: Setup,
    Verify, Sign, Collect, Finalize, Verify.  Each signer deposits tokens
    independently, and the authority finalizes atomically.
  \item The eUTxO concurrency problem---where multiple signers competing for
    the same UTxO causes all but one to fail---is solved by having each
    signer create an independent UTxO at the validator address.
  \item The \texttt{SignerDatum} structure records five fields (signer PKH,
    document hash, policy ID, validity expiry, deposit slot) that together
    prove who signed, what they signed, when, and with what credentials.
  \item Privacy is achieved through strict architectural separation: only
    pseudonymous hashes and policy IDs go on-chain; all personal data stays
    in an AES-256-GCM encrypted off-chain database.
  \item GDPR compliance is achieved through crypto-shredding: destroying
    per-credential encryption keys renders off-chain personal data
    permanently inaccessible while on-chain hashes remain meaningless
    without the mapping.
  \item The ESIGN Act requires intent, consent, association, and retention---all
    satisfied by the system's deliberate signing process, onboarding flow,
    hash binding, and blockchain immutability.
  \item UETA, adopted by 49 states plus DC, provides the state-level legal
    foundation; nine states have blockchain-specific amendments with Arizona
    and Vermont being the most significant.
  \item eIDAS~2.0 defines three signature tiers (SES, AES, QES).  The
    credential system natively meets all four AES requirements.  Article~45l
    provides a pathway to qualified trust service status through the new
    ``Qualified Electronic Ledger'' category.
\end{itemize}
